\documentclass{article}
\usepackage{amsmath, graphicx}
\usepackage{grffile}
\usepackage{dcolumn}
\usepackage{bm}
\usepackage{url}
\usepackage{epsfig}
\usepackage{mathrsfs}  %  package for the "curly" fonts
\usepackage{subfigure}
\usepackage{multirow}
\usepackage{algorithm}
\usepackage{algorithmicx}
\usepackage{amssymb}
\usepackage{tensor}
\usepackage{romannum}

\usepackage{amsfonts}
\usepackage{amssymb}
\usepackage{graphicx}
\usepackage{algpseudocode}
\usepackage{algorithm}
\usepackage{algorithmicx}
\usepackage{color}
\usepackage[left=2cm,right=2cm,top=2cm,bottom=2cm]{geometry}
\usepackage{fancyhdr}
\usepackage{subfiles} 

\usepackage{hyperref}
\hypersetup{
    colorlinks=true,
    linkcolor=blue,
    filecolor=magenta,      
    urlcolor=cyan,
}

\usepackage{listings}
\lstset{language=Python, numbers=left, numberstyle=\tiny, stepnumber=1, numbersep=5pt, tabsize=4, basicstyle=\ttfamily}
\usepackage{inconsolata}


\topmargin 0.30in
\textheight 9.00in
\parindent 0.0in

 \addtolength{\voffset}{-2cm}
 
%Some macro
% A useful macro for TODOS
\newcommand{\TODO}[1]{{\color{red}[}{\color{red}TODO:} {\color{blue}#1}{\color{red}]}}
% A note we can use to discuss the draft with each other in the pdf
\newcommand{\NOTE}[1]{{\color{blue}[#1]}}
% Convenient macros for autoscaling containers
\newcommand{\paren}[1]{\left(#1\right)}
\newcommand{\sqrbrc}[1]{\left[#1\right]}
% inverse
\newcommand{\inv}[1]{\frac{1}{#1}}
% derivatives
\newcommand{\dd}[1]{\frac{d}{d#1}}
\newcommand{\pdd}[1]{\frac{\partial}{\partial #1}}
% use this for vectors so we can swap formatting easily
\newcommand{\myvec}[1]{\vec{#1}}
% Big O notation
\newcommand{\Ord}[1]{\mathcal{O}\paren{#1}}


\begin{document}

\title{WVT-like relaxation in FleCSPH}
\author{Irina Sagert}

\maketitle

Choosing optimal initial conditions for SPH in terms of particle placement is important for e.g. noise reduction and minimization of numerical artifacts. 
To avoid geometrical effects when particles are placed on a regular lattice (for example a dependence of results on grid orientations), lattice positions are often perturbed, followed by a relaxation stage where particles find a new energetically minimal state. \\
FleCSPH currently has two relaxation techniques. 
In one, perturbed particles are confined in an external potential which is set up in such a way that the particles are pushed towards a target density profile.  
Another option that is discussed here is the relaxation via direct repulsive interaction of particles that lie within a smoothing length of one another. 
The repulsive forces can have different forms but are often functions of smoothing lengths and particle distance. 
This method has been shown to lead to fast relaxation where the final particle positions are similar to a  Weighted Voronoi Tessellation placement, hence terming it WVT-like.
In our implementation we follow the work of Diehl et al. \cite{Diehl15} and Arth et al. \cite{Arth19} and discuss details below. 
%%%%%%%%%%%%%%%%%%%%%%%%%%%%%
%%%%%%%%%%%%%%%%%%%%%%%%%%%%%
\section{Repulsive inter-particle forces}
%%%%%%%%%%%%%%%%%%%%%%%%%%%%%
%%%%%%%%%%%%%%%%%%%%%%%%%%%%%
\subsection{Diehl et al.}
%%%%%%%%%%%%%%%%%%%%%%%%%%%%%
For a fast relaxation process, Diehl et al. \cite{Diehl15} introduced a repulsive inter-particle force with a $r^{-2}$ dependence, similar to a repulsive gravitational force. At each iteration, the displacement of particle $i$ is given by:
%%%%%%%%%%%%%%%%%%%%%%%%%%%%%
\begin{align}
\Delta \mathbf{r}_i = \mu_\mathrm{wvt}  \: D_i, \:\:\: D_i = h_i \: \sum_j \left[\left(\frac{h_{ij}}{r_{ij} + \epsilon}\right)^2 + \kappa \right] \frac{\mathbf{r}_{ij}}{|\mathbf{r}_{ij}|}, \:\:\: \kappa = - \left(\frac{h_{ij}}{h_{ij} + \epsilon}\right)^2 
\label{diehl_deltax}
\end{align}
%%%%%%%%%%%%%%%%%%%%%%%%%%%%%
where $h_{ij} = (h_i + h_j)/2$, $\mathbf{r}_{ij} = \mathbf{r}_i - \mathbf{r}_j$. The sum in eq.(\ref{diehl_deltax}) runs over all neighbors $j$. The factor $\mu_\mathrm{wvt}$ specifies the fraction of $h_i$ that a particle $i$ is allowed to move during each iteration step. 
The value of $\epsilon$ can be chosen as a small fraction of $h_{ij}$ to avoid numerical problems for particles that are close together. 
The constant $\kappa$ is added to the force to ensure that the latter vanishes for $| \mathbf{r}_{ij}| > h_{ij}$.  
Diehl et al. recommend to gradually decrease the value of $\mu_\mathrm{wvt}$ with each iteration, allowing large position displacements in the beginning of the relaxation process and restricting the particle motion as the relaxed state is approached.  
For simplicity, we leave the value of $\mu_\mathrm{wvt}$ a constant but set it to a small value. 
An option for linearly decreasing  $\mu_\mathrm{wvt}$ once the simulations is close to convergence is implemented via the \texttt{wvt\_cool\_down} iterations were we determine a new $\mu_\mathrm{wvt}$ via:
%%%%%%%%%%%%%%%%%%%%%%%%%%%%%
\begin{align}
\mu_\mathrm{wvt, new} = \mathrm{max}[a \times \mathrm{iterations} + b,0.0], \:\:\: a = - \mu_\mathrm{wvt}/(\mathrm{wvt\_cool\_down}), \:\:\: b = \mu_\mathrm{wvt}.
\end{align}
%%%%%%%%%%%%%%%%%%%%%%%%%%%%%
However, tests have not yet shown that this leads to a better convergence or improved final particle configuration. 
At present, we leave it as an option for the user. 
%%%%%%%%%%%%%%%%%%%%%%%%%%%%%
\subsection{Arth et al.}
%%%%%%%%%%%%%%%%%%%%%%%%%%%%%
The WVT-like relaxation method introduced in Arth et al. \cite{Arth19} is very similar to \cite{Diehl15}. 
The largest difference lies in the form of the repulsive inter-particle force as Arth et al. use the SPH kernel itself. 
At each iteration, the particle displacement is calculated via:
%%%%%%%%%%%%%%%%%%%%%%%%%%%%%
\begin{align}
\Delta \mathbf{r}_i = \mu_\mathrm{wvt} \: A_i, \:\:\: A_i = \sum_j h_{ij} \: W(|\mathbf{r}_{ij}|, h_{ij})  \: \frac{\mathbf{r}_{ij}}{|\mathbf{r}_{ij}|}.
\label{arth_deltax}
\end{align}
%%%%%%%%%%%%%%%%%%%%%%%%%%%%%
where the expressions for $h_{ij}$, $\mathbf{r}_{ij}$ and $ \mu_\mathrm{wvt}$ are as for Diehl et al.
We implement both methods for the repulsive particle force. 
The corresponding input deck parameter is \texttt{wvt\_method} which can be either set to \texttt{diehl} or \texttt{arth}. 
Note that because of the different formulations of the repulsive inter-particle force, the two methods need different values for the wvt input deck parameters which we will discuss below. 
A parameter set that works well for the method of Arth et al. is not guaranteed to work as well for the repulsive force of Diehl et al. and most likely needs to be re-adjusted. 
%%%%%%%%%%%%%%%%%%%%%%%%%%%%%
%%%%%%%%%%%%%%%%%%%%%%%%%%%%%
\section{Normalization of the Smoothing Length}
%%%%%%%%%%%%%%%%%%%%%%%%%%%%%
%%%%%%%%%%%%%%%%%%%%%%%%%%%%%
According to Diehlt et al. and Arth et al., the value of the smoothing length should be normalized to a given number of SPH particles $N_\mathrm{total}$ and a desired number of neighbors $N_\mathrm{ngb}$.
The latter is set via the parameter \texttt{wvt\_ngb}. 
At each iteration, we determine the target density for particle locations $\rho_{t,i} = \rho_t (\mathbf{r}_i)$ and the corresponding smoothing lengths. 
As is done in Diehl et al. \cite{Diehl15} and in our regular SPH formalism we calculate the latter via:
\begin{align}
h_{t,i} = \eta \left(\frac{m}{\rho_{t,i}}\right)^{1/3}.
\label{h_diehl}
\end{align}
However, Arth et al. uses a different formulation that is tied to the desired number of neighbors:
%%%%%%%%%%%%%%%%%%%%%%%%%%%%%
\begin{align}
h_{t,i} = \left( \frac{N_\mathrm{nbg} \: m}{\frac{4}{3} \pi \rho_{t,i}} \right)^{1/3} .
\label{h_arth}
\end{align}
%%%%%%%%%%%%%%%%%%%%%%%%%%%%%
We choose eq.(\ref{h_diehl}) as the default expression but include eq.(\ref{h_arth}) as an option. 
The corresponding parameter is \texttt{wvt\_h\_ngb} which should be set to \texttt{true} to determine the smoothing lenght via eq.(\ref{h_arth}). 
After the target smoothing lengths have been set the sum of all individual SPH particle volumes is calculated via
%%%%%%%%%%%%%%%%%%%%%%%%%%%%%
\begin{align}
V_\mathrm{SPH} = \frac{4}{3} \pi \: \sum_i h_{t,i}^3 .
\end{align}
%%%%%%%%%%%%%%%%%%%%%%%%%%%%%
The smoothing lengths for all particles are then scaled so that $V_\mathrm{total} = V_\mathrm{SPH} / N_\mathrm{ngb}$ which means that they are multiplied by a scaling factor $\alpha$.
With $V_\mathrm{total} = (4 \pi/3) \: r_\mathrm{total}^3$ the scaling factor is:
%%%%%%%%%%%%%%%%%%%%%%%%%%%%%
\begin{align}
h_i = \alpha \: h_{t,i}, \:\:\: \alpha =  \left( \frac{N_\mathrm{ngb}}{\sum_i h_{t,i}^3} \right)^{1/3} r_\mathrm{total}.
\end{align}
%%%%%%%%%%%%%%%%%%%%%%%%%%%%%
%%%%%%%%%%%%%%%%%%%%%%%%%%%%%
\section{Boundary Conditions}
%%%%%%%%%%%%%%%%%%%%%%%%%%%%%
%%%%%%%%%%%%%%%%%%%%%%%%%%%%%
Using repulsive particle interactions for relaxation requires the usage of boundary conditions in order to confine the SPH particles within the simulated object. 
Diehl et al. recommend to use ghost particles as these are a common way to implement boundary conditions in SPH codes. 
The ghost particles can be created by mirroring simulation particles that lie within a distance of one or two smoothing lengths from the boundary. 
Arth et al. choose a different approach and apply periodic boundary conditions for relaxation cutting out the region of interest once convergence is reached. \\
%%%%%%%%%%%%%%%%%%%%%%%%%%%%%
For the WVT relaxation in FleCSPH, we currently have two boundary options.  
They can be set by the parameter \texttt{wvt\_boundary} to either \texttt{reflective} or \texttt{frozen} and currently are only guaranteed to work for spherical objects, e.g. neutron stars, white dwarfs, simple spheres or disks.
In the \texttt{reflective} setting, we mirror the neighbors of each particle with $r_a+h_a \leq r_\mathrm{sphere}$ at the sphere boundary, where $r_a$ is the radial distance of the particle, $h_a$ its smoothing length, and $r_\mathrm{sphere}$ the radius of the sphere. 
The particle of interest can interact with the mirror images of its neighbors as well as itself if the mirror images fulfill the same conditions as the regular neighbor particles.\\
Spherical mirroring at a sphere's radius as implemented for the \texttt{reflective} boundary conditions is straight forward but leads to a slightly lower density of the mirrored particles. 
This could affect the relaxation process or lead to a deviation of the relaxed particle density in comparison to the target density profile, especially close to the spherical boundary. 
An option that guarantees the same density of the reflected particles is to usage of a plane mirror that is normal to a particles position vector. 
Such a mirror could be positioned either at the intersection of the particle position vector with the sphere or at the location of the particle. 
Due to its simple implementation, we chose the latter. 
Repulsive forces of particles and mirror particles that are parallel to a particle's position vector cancel out and only contributions that are normal to the position vector remain. 
We implement this boundary condition as \texttt{frozen} by computing the total displacement for a particle due to its neighbors, removing the contribution that is parallel to the particle position vector and doubling the one that is normal.  
%%%%%%%%%%%%%%%%%%%%%%%%%%%%%
%%%%%%%%%%%%%%%%%%%%%%%%%%%%%
\section{Convergence Check}
%%%%%%%%%%%%%%%%%%%%%%%%%%%%%
%%%%%%%%%%%%%%%%%%%%%%%%%%%%%
To monitor the relaxation process and implement a condition for its termination, we apply a convergence check from Arth et al. \cite{Arth19}. 
Instead of e.g. using the error in the density, the authors suggest to stop the relaxation when most particles move less than a small fraction of the mean particle spacing $d_{\mathrm{mps},i}$. 
The latter is calculated via:
%%%%%%%%%%%%%%%%%%%%%%%%%%%%%
\begin{align}
	d_{\mathrm{mps},i} = \left(\frac{1}{N_\mathrm{ngb}}\right)^{1/3} \: h_i. 
\end{align}
%%%%%%%%%%%%%%%%%%%%%%%%%%%%%
In the code, we count the number of particles $N_\mathrm{cnt}$ that experience a displacement $|\Delta \mathbf{r}_i | >  10^{-3} \: d_{\mathrm{mps},i}$ in each iteration. The number is expressed as a percentage of the total particle number: 
%%%%%%%%%%%%%%%%%%%%%%%%%%%%%
\begin{align}
	m_\mathrm{mps} = 100 \: (N_\mathrm{cnt} / N_\mathrm{total}). 
\end{align}
%%%%%%%%%%%%%%%%%%%%%%%%%%%%%
When $m_\mathrm{mps}$ is smaller than a user-specified value, the relaxation process is stopped. 
The convergence check is switched on by setting the parameter \texttt{wvt\_convergence\_check} to \texttt{true}. 
The convergence percentage is given by the value of \texttt{wvt\_convergence\_point}. 
Values that are around one seem to work well. 
If the above parameters are not set or \texttt{wvt\_convergence\_check} is \texttt{false}, the simulation will simply run for the number of iterations that are given by \texttt{final\_iteration}.
%%%%%%%%%%%%%%%%%%%%%%%%%%%%%
%%%%%%%%%%%%%%%%%%%%%%%%%%%%%
\section{Code Implementation}
%%%%%%%%%%%%%%%%%%%%%%%%%%%%%
%%%%%%%%%%%%%%%%%%%%%%%%%%%%%
The WVT-like relaxation is implemented into FleCSPH using the regular code architecture in the wvt driver.  
The detailed algorithm is given below and 
%%%%%%%%%%%%%%%%%%%%%%%%%%%%%
\begin{algorithm}
\caption{WVT relaxation}
\label{alg:wvt}
\begin{algorithmic}[1]
\State read initial state: {\tt \{r, v, u\}} $\leftarrow \{r^0, v^0, u^0\}$
\State compute density {\tt rho(r)} for output only
\State compute $D_i$ or $A_i$
\While{iterations $<$ final\_iterations}
\State compute displacement $|\Delta \mathbf{r}_i|$; move particles with boundary conditions; set target {\tt rho\_i(r)} and {\tt h\_i} 
\State rescale smoothing length 
\State compute density {\tt rho(r)} for output 
\State compute $D_i$ or $A_i$
\State check convergence; set iterations to final\_iterations if convergence check is fulfilled
\EndWhile
\While{iterations $<$ (final\_iterations + wvt\_cool\_down)}
\State compute displacement $|\Delta \mathbf{r}_i|$; move particles with boundary conditions; set target {\tt rho\_i(r)} and {\tt h\_i} 
\State rescale smoothing length 
\State compute density {\tt rho(r)} for output only
\State compute $D_i$ or $A_i$
\EndWhile
\end{algorithmic}
\end{algorithm}
%%%%%%%%%%%%%%%%%%%%%%%%%%%%%
it usually takes about $100 < \mathrm{iterations} < 2000$ for the relaxation process. 
Note that although we take advantage of the FleCSPH code structure and typical SPH quantities like smoothing lengths and neighbors to calculate repulsive forces, the WVT-like relaxation is not a SPH simulation! 
This means e.g. that the density as reported in steps 2, 7 and 14 does not necessarily correspond to the true density in the final SPH simulation. 
It is reported for output purposes only and only to give the user an impression of the relaxation process. 
The user should always run a modification script followed by a few iterations steps with the hydro or newtonian driver to examine the final density. 
An example of input, relaxation and modification input decks is given in the next section. 
%%%%%%%%%%%%%%%%%%%%%%%%%%%%%
%%%%%%%%%%%%%%%%%%%%%%%%%%%%%
\section{Example Input Files}
%%%%%%%%%%%%%%%%%%%%%%%%%%%%%
%%%%%%%%%%%%%%%%%%%%%%%%%%%%%
The following input file creates a disk (2D) or sphere (3D) with the mesa density profile. 
It can be initialized using the Sedov test generator and run with the wvt driver. 
%%%%%%%%%%%%%%%%%%%%%%%%%%%%%
\begin{lstlisting}
# ##############################################
# Input deck to create a 3D sphere / 2D disk with a density 
# that follows the mesa radial density profile. 
# Particles are initialized with random positions followed by
# WVT relaxation 
#
# Usage: ./sedov_3d_generator <this_file>.par for initial data
#        ./wvt_3d <this_file>.par for wvt relaxation
##############################################

# initial data
  initial_data_prefix = "wvt_initial"
  initial_iteration = 0

# geometry: spherical domain with radius R = 1
  domain_type = 1                           # 0:box, 1:sphere
  sphere_radius = 1.0                       # radius of sphere 
  density_profile = "mesa"
  external_force_type = "spherical density support"

# icosahedral lattice with small perturbations
  lattice_nx = 60                           # 60x60x60 particles (3D)
  lattice_type = 4                          # 0:rect.; 1:hcp; 2:fcc; 3:ico; 4:rand

# equation of state type and parameters
  eos_type = "no eos"                       # no eos needed for wvt 

# density and pressure
  rho_initial = 1.0                         # desired central density
  pressure_initial = 0.1                    # desired central pressure 

# set Sedov energy to zero
  sedov_blast_energy = 0.0

# parameters for relaxation / evolution
  final_iteration = 2000                     # number of iterations 
  
  wvt_method = "arth"                        # "arth" or "diehl"
  wvt_boundary = "reflective"                # "reflective" or "frozen"
  wvt_mu = 1.e-3                             # small fraction of smoothing length 
  wvt_ngb = 140                              # desired number of neighbors within h 
  wvt_convergence_check = true
  wvt_convergence_point = 1.0
  wvt_cool_down = 0
  wvt_h_nbg = false

# output parameters
  out_screen_every = 20                     # output on screen
  out_scalar_every = 20                     # output in .dat files
  out_h5data_every = 20                     # output in h5part files
  output_h5data_prefix = "wvt_3d"

# SPH kernel parameters
  sph_kernel = "Wendland C6"                # kernel choice
  sph_eta = 1.6                             # smoothing length
  sph_variable_h = true                     # variable smoothing length
\end{lstlisting}
%%%%%%%%%%%%%%%%%%%%%%%%%%%%%
After running the initialization and relaxation steps it is recommended to run a modification script to reset the values of the particle specific internal energy according to the desired equation of state. An example of a modification script is:
%%%%%%%%%%%%%%%%%%%%%%%%%%%%%
\begin{lstlisting}
# ##############################################
# Modification after WVT relaxation
#
# Usage: ./sedov_3d_generator <this_file>.par to 
# modify data
# ##############################################

  initial_data_prefix = "wvt_3d"
  modify_initial_data = yes                 # flag to modify data
  initial_iteration = 360                   # set to final output iteration of wvt

# geometry: spherical domain with radius R = 1
  domain_type = 1                           # 0:box, 1:sphere
  sphere_radius = 1.0                       # radius of sphere 
  density_profile = "mesa"
  external_force_type = "spherical density support"

# equation of state type and parameters
  poly_gamma = 1.4                          # eos parameter

# density and pressure
  rho_initial = 1.0                         # desired density
  pressure_initial = 0.1                    # desired pressure 

# since we only need spherical distribution of particles,
# set Sedov energy to zero
  sedov_blast_energy = 0.0

# parameters for relaxation / evolution
  final_iteration = 20                      # number of interations 

# output parameters
  out_screen_every = 20                     # output on screen
  out_scalar_every = 20                     # output in .dat files
  out_h5data_every = 1                      # output in h5part files
  output_h5data_prefix = "wvt_ev"

# SPH kernel parameters
  sph_kernel = "Wendland C6"                # kernel choice
  sph_eta = 1.6                             # smoothing length
  sph_variable_h = yes                      # variable smoothing length
\end{lstlisting}
%%%%%%%%%%%%%%%%%%%%%%%%%%%%%
%%%%%%%%%%%%%%%%%%%%%%%%%%%%%
%%%%%%%%%%%%%%%%%%%%%%%%%%%%%
\section{Problems and Solutions}
%%%%%%%%%%%%%%%%%%%%%%%%%%%%%
%%%%%%%%%%%%%%%%%%%%%%%%%%%%%
The WVT-like relaxation is still experimental and  does not work perfectly in the sense that the particle densities do not relax exactly to the target density profile but scatter around it.
A relaxation method that gives final states with very little noise is the relaxation in an external potential. 
However, the WVT-like method works well for stellar profiles, especially with a large density gradients which can be challenging for the external potential relaxation. Convergence, as previously defined for WVT, is usually reached within less than 2000 iterations. \\
However, some known problems with WVT in FleCSPH are listed below, together with possible solutions:
\begin{itemize}
\item \textbf{Problem}: The particle densities of the relaxed configuration scatter around the target densities, but the noise is too large. \textbf{Solution}: This usually occurs for too few iteration steps, particles or neighbors. The user can first try to decrease the value of \texttt{wvt\_convergence\_point} and see if this leads to an improvement. As a next step, the number of \texttt{wvt\_ngb} can be increased. Note that these both of these steps have been seen to reduce the noise in density, but, for density profiles of stars, e.g. white dwarfs, also lead to a larger mismatch in the density profile with respect to the target densities. If this is observed and the mismatch is too large to be acceptable, the only solution that is know to work at present is to increase the total particle number
\item \textbf{Problem}: The relaxed density profile does not match the target densities well; in particular, the relaxed particle densities are lower than the target densities. \textbf{Solution}: This has been seen for stellar density profiles with large density gradients, e.g. white dwarfs, when too many iterations steps are performed and/or too many wvt neighbors are used. An increase in \texttt{wvt\_mu} or decrease of \texttt{wvt\_ngb} can help. Alternatively, the user can set \texttt{wvt\_convergence\_check} to \texttt{false}, relax a configuration for ca. 2000 iterations and select the density profile that the users sees best fit. 
\item \textbf{Problem}: The \texttt{wvt\_convergence\_check} is switched on but the simulation does not converge. \textbf{Solution}: This has been seen for some cases when the number of wvt neighbors was set to a high value. A decrease of \texttt{wvt\_ngb} should be attempted.  
\end{itemize}


\bibliographystyle{unsrt}
\bibliography{wvt_refs.bib}

\end{document}
